\clearpage
\begin{summarybox}
\textbf{一分钟记住：}
\[
\boxed{\ f(x+a)=-f(x)\quad\Longrightarrow\quad f(x+2a)=f(x).\ }
\]
\textbf{口诀：走一步变号，走两步还原；偶数步不变，奇数步相反。}

\medskip
\textbf{三个最常用的判断：}
\begin{enumerate}[leftmargin=*,itemsep=0.3em]
\item \textbf{求周期：} $2|a|$ 是一个正周期，\textbf{未必是最小正周期}。
\item \textbf{求函数值：} 输入相差 $na$，若 $n$ 为偶数就等于原值；若 $n$ 为奇数就取原值的相反数。
\item \textbf{求一长串和：} $f(x)+f(x+a)=0$，两项一组往往马上消掉。
\end{enumerate}

\medskip
\textbf{一个关键防错：}只凭变号条件，\textbf{不能}判断图像关于 $\bigl(\frac a2,0\bigr)$ 中心对称。若另知 $f$ 为\textbf{偶函数}，才得到该对称中心；若另知 $f$ 为\textbf{奇函数}，得到对称轴 $x=\frac a2$。

\medskip
\textbf{考试写法（最短两行）：}
\[
\begin{aligned}
f(x+2a)&=-f(x+a)\\
&=-[-f(x)]=f(x).
\end{aligned}
\]
因此 $2|a|$ 是函数的一个正周期（$a\ne0$）。

\medskip
\textbf{错位反号速查：}
\[f(x+a)=-f(x+b),\ a\ne b
\quad\Longrightarrow\quad \boxed{2|a-b|\text{ 是一个正周期。}}\]
先比较两输入的距离，再走两次；以上均要求关系对全部输入成立、定义域允许所需的来回平移，且所得不保证最小。
\end{summarybox}
